% New, standalone full manuscript. Original manuscript files remain untouched.
% The empirical section distinguishes archived prototype runs from the revised
% mechanisms, whose implementation and matched evaluation are still required.
\documentclass[10pt,twocolumn]{article}
\usepackage[letterpaper,margin=0.75in,columnsep=0.25in]{geometry}
\usepackage[T1]{fontenc}
\usepackage{lmodern}
\usepackage{amsmath,amssymb,amsthm}
\usepackage{booktabs,array}
\usepackage{framed}
\usepackage[hidelinks]{hyperref}
\setlength{\emergencystretch}{2em}
\setlength{\parindent}{1em}
\setlength{\parskip}{2pt}
\newcommand{\sys}{MemCatalog}
\newcommand{\Hh}{\mathcal{H}}
\newcommand{\Ff}{\mathcal{F}}
\newcommand{\Uu}{\mathcal{U}}
\newcommand{\Mm}{\mathcal{M}}
\newcommand{\ret}{\textsc{Retrieve}}
\newcommand{\red}{\textsc{Reduce}}
\newcommand{\comp}{\textsc{Compose}}
\newcommand{\runin}[1]{\smallskip\noindent\textbf{#1.}\ }
\DeclareMathOperator*{\argmin}{arg\,min}
\DeclareMathOperator*{\argmax}{arg\,max}
\newtheorem{proposition}{Proposition}
\newtheorem{definition}{Definition}
\newenvironment{memoryexample}[1]{%
  \def\FrameCommand{\fboxrule=0.7pt\fboxsep=6pt\fbox}%
  \begin{framed}\small\noindent\textbf{#1}\par\smallskip
}{\end{framed}}

\title{\sys: Query-Driven Facts and Second-Order Memory\\
for LLM Agents}
\author{Anonymous Author(s)}
\date{}

\begin{document}
\maketitle

\begin{abstract}
An agent's conversation history contains more information than any fixed
collection of extracted memories can anticipate. Moreover, knowing facts
about a user does not reveal which facts that user's agent needs, how they
should be combined, or how much computation a memory query deserves.
\sys{} addresses these problems with three coupled forms of persistent
state: raw sessions preserve the evidence, first-order memory acquires
reusable facts as queries reveal a need for them, and second-order memory
remembers how those facts were retrieved and used. A dependency-aware
planner expresses memory queries through evidence retrieval, reduction,
and composition, each with an explicit execution pipeline. An
experience-guided executor allocates computation to individual queries
and revises that allocation as evidence and uncertainty change. Each
interaction therefore improves both the agent's knowledge of the user
and the memory system's knowledge of how to serve that user. We specify
the architecture, the three recursive pipelines, and a sequential
risk--cost decision rule. An audit of 2,261 archived records establishes
the empirical starting point of the released prototype; these records
are distinguished from the evaluation still required for the revised
fact-use and budget mechanisms.
\end{abstract}

\section{Introduction}
\label{sec:new-introduction}

An LLM agent serving a user over many conversations needs to recover
facts that its current context no longer contains. Those facts range
from a single preference to a history of activities, changing plans,
and relationships between people and events. Long-term conversation
benchmarks show that recalling such information involves information
extraction, temporal reasoning, and reasoning across sessions, rather
than one uniform retrieval task~\cite{locomo,longmemeval}. The central
problem is how a memory system acquires the right facts, combines them
for a particular query, and improves its service as the user continues
interacting with the agent.

Existing systems provide useful mechanisms for this problem. MemGPT
manages information across memory tiers~\cite{memgpt}; Mem0 extracts and
updates salient facts~\cite{mem0}; Zep preserves temporal relationships
in a knowledge graph~\cite{zep}; and A-MEM develops linked memory
notes~\cite{amem}. Retrieval controllers such as MemR$^3$ explicitly
refine retrieval and track evidence gaps~\cite{memr3}. These mechanisms
establish that memory construction and access can be adaptive. Our
question is how to connect \emph{query-driven fact acquisition},
\emph{experience of using those facts}, and \emph{allocation of
computation} in one explicit memory interface.

\runin{The next useful fact cannot always be predicted at write time}
A conversation can mention an activity's organizer, location, duration,
and the user's motivation. An early query may require only its
organizer; a later query may ask which activity lasted longest. A fixed
summary can omit the duration, while extracting every possible fact
upfront spends effort on information that may never be requested.
The original history must remain available, and incoming queries must
be able to reveal new factual requirements. Query-driven evolution,
as articulated in EUR~\cite{eur}, suggests the useful principle that
a query can both consume existing knowledge and leave behind new
reusable knowledge. For agent memory, this principle concerns the
user facts an agent needs over time.

\runin{Facts do not remember how they were useful}
A fact store describes the user, but does not by itself describe the
user's memory workload. A frequently retrieved event might never
contribute to an answer, whereas an otherwise obscure organizer fact
might repeatedly connect relevant activities. Retaining this distinction
requires a second form of memory: experience of which facts were
exposed, which were actually used, and under what information need.
Such experience can help prioritize future evidence work. It must be
kept separate from the truth of a fact, because frequency of use does
not establish source support.

\runin{The required computation changes with the query}
``Who organized my hike?'' selects a fact. ``How many outdoor events
did I attend?'' must cover a population and decide a condition for each
candidate. ``Who organized the longest outdoor event?'' must identify
the event before retrieving its organizer. A single prompt over
retrieved passages hides these different obligations. The memory
system needs an explicit plan that reveals which facts to obtain,
which values to reduce, and which intermediate results to connect.

\runin{A plan label does not determine execution difficulty}
Comparing two cached durations needs little work. Resolving a name
shared by two people may require additional evidence even though it
looks like a direct retrieval. A useful execution budget therefore
depends on the individual query, its current evidence, and previous
experience, then changes as the execution resolves or reveals
uncertainty. Repeated samples are justified by their expected benefit,
not by assigning one constant to every query in a class.

We propose \sys{}, a memory management system built around these
observations. Its state has three coupled components. Raw sessions
preserve the conversational evidence; first-order memory acquires facts
as queries require them; and second-order memory represents experience
of acquiring, retrieving, and using those facts. The agent poses a
memory query, receives a supported result and its evidence, and reports
which facts contributed to its answer. This closes the loop from an
information need to lasting knowledge and use experience.

The architecture supports two technical mechanisms. For performance,
queries are represented through evidence retrieval, reduction, and
composition. Their recursive structure defines the required execution
dependencies; contextual fact-use experience determines where evidence
work begins. For efficiency, an executor predicts evidence risk and
execution cost, allocates a query-specific budget, and spends additional
work only where the predicted improvement warrants its cost. The
controller reuses verified subresults and targets unresolved blocks.

The contributions are: (1) a three-part memory model and an explicit
agent--memory interface that make query-driven facts and second-order
use experience complementary forms of persistent state; (2) a recursive
plan language with disjoint root classes, three specified evidence
pipelines, and a precise structural coverage claim; and (3) an
experience-conditioned budget objective and a sequential stopping rule
with a local risk--cost justification. Section~\ref{sec:new-experiments}
audits the available prototype records and defines the matched
evaluation needed to test the revised mechanisms. The archived runs
do not by themselves establish their performance or cost benefits.

\section{Preliminaries and Problem Formulation}
\label{sec:new-preliminaries}

\runin{Scope and source availability}
We consider one agent serving one user over an evolving sequence of
timestamped conversation sessions. A memory query can use only sources
available when it arrives. Its requested temporal scope can be current,
historical, or a specified interval. Facts stated later are excluded
even if they would clarify an earlier event. External retrieval and
creative generation can be agent tools, but are outside the memory
query model studied here.

\runin{Tasks and evidence obligations}
A memory task specifies an information request and the evidence needed
to support it. Direct requests need a supported fact and identity;
population requests additionally need candidate coverage; relational
requests need supported connections among subresults. A task can be
unanswerable because the history lacks the evidence or because a finite
execution budget has not recovered it. A response exposes these limits
through its support status instead of assuming a closed-world fact store.

\runin{Objective}
For the $t$th question, the agent generates $A_t$ from $Q_t$ and the
retrieved memories. With a held-out answer-quality function
$\operatorname{Qual}$ and memory cost $C_t$, the design objective is
\begin{equation}
\label{eq:new-objective}
\max\ \mathbb{E}\!\left[\sum_t
       \bigl(\operatorname{Qual}(Q_t,A_t)-\lambda C_t\bigr)\right].
\end{equation}
Quality is measured against reference answers in evaluation; those
answers are unavailable to the deployed system. The executor therefore
uses observable evidence-risk diagnostics as a surrogate, as defined
in Section~\ref{sec:new-efficiency}. Its predictions can be wrong and
must be assessed against actual answer quality. Per-user experience
can improve service, but monotone improvement is not an assumption of
the model.

\begin{table}[t]
\centering\small
\caption{The core notation. Uppercase user messages and lowercase memory
messages distinguish the two interfaces.}
\label{tab:new-notation}
\begin{tabular}{@{}ll@{}}
\toprule
Symbol & Meaning\\
\midrule
$Q,A$ & User question and agent answer\\
$q,m$ & Memory query and retrieved memory\\
$\Hh,\Ff,\Uu$ & Sessions, facts, use experience\\
$P$ & Recursive memory query plan\\
$E^r,E^u$ & Retrieved and used fact sets\\
$\ell,C$ & Evidence-risk diagnostic and cost\\
$K(q)$ & Query-specific execution budget\\
\bottomrule
\end{tabular}
\end{table}

\runin{Organization}
Architecture defines the interface and the state transition.
Performance specifies the plan language, second-order grounding,
and three execution pipelines. Efficiency controls how much work those
pipelines receive. The notation in Table~\ref{tab:new-notation} is reused
throughout; facts, events, and attributes do not require separate
families of sets.

\section{Architecture}
\label{sec:new-architecture}

\sys{} is a memory management system for an LLM agent. Its architecture
separates three things that evolve at different rates: what the user has
said, which facts the agent has needed, and how those facts have helped
the agent. This section defines the interface and the persistent state;
Sections~\ref{sec:new-performance} and~\ref{sec:new-efficiency} explain
how a query changes that state and how its execution is budgeted.

\subsection{The Agent--Memory Interface}
\label{subsec:new-interface}

\runin{Question and answer; query and retrieved memory}
We denote a user-facing \emph{question} and \emph{answer} by $Q$ and $A$,
and an agent-facing \emph{memory query} and \emph{retrieved memory} by
$q$ and $m$. A question expresses what the user wants; a memory query
specifies what information the agent needs to answer it. The two may
coincide for a factual question, but a recommendation can induce several
queries about preferences, constraints, and past activities. The
terminology follows the distinction between a retrieval query and the
memories or evidence returned by a memory system~\cite{mem0,memr3}.
The interaction is
\begin{equation}
\label{eq:new-interface}
\begin{aligned}
\mathrm{User}&\xrightarrow{\ Q\ }\mathrm{Agent}
                    \xrightarrow{\ q\ }\Mm,\\
\Mm&\xrightarrow{\ m\ }\mathrm{Agent}
                    \xrightarrow{\ A\ }\mathrm{User}.
\end{aligned}
\end{equation}
The agent issues queries and generates the final answer. The memory
system supplies facts, computed results when requested, and their
supporting evidence. It also receives a report of which supplied facts
the agent used. The relevant user--agent exchange is subsequently stored
as part of a raw session.

\runin{A memory query has an evidence scope}
The query $q$ contains an information request, an entity or event scope,
and a temporal scope. Unless specified otherwise, its evidence is drawn
from the user's history available when the query arrives. ``Where did
I live in January?'' and ``Where do I live now?'' can refer to the same
attribute while selecting different fact versions. A query requiring
exhaustive coverage also identifies the population over which a result
must be computed, such as all events in March.

\runin{Retrieved memory is a supported result}
We represent the response as
\begin{equation}
\label{eq:new-response}
m=(v,E,\sigma),
\end{equation}
where $v$ is the requested value or evidence collection, $E\subseteq\Ff$
contains the supporting first-order facts, and $\sigma$ indicates
\texttt{supported}, \texttt{partial}, or \texttt{unknown}. Facts in $E$
retain links to their source spans. A count is returned with its
contributing events; a comparison with the values compared. When a
requirement is unresolved, the response identifies the gap rather than
turning a missing fact into a negative assertion. ``No event matched''
requires coverage of the requested population; ``no matching event has
yet been found'' does not establish that claim.

\subsection{Three Coupled Forms of Memory}
\label{subsec:new-state}

\runin{The state}
For one user, the persistent memory state is
\begin{equation}
\label{eq:new-state}
\boxed{\Mm=(\Hh,\Ff,\Uu).}
\end{equation}
The paired roles are simple: $\Hh$ preserves \emph{conversations}, $\Ff$
stores \emph{facts about the user}, and $\Uu$ stores \emph{experience of
using those facts}. All state and updates are isolated by user. We omit
the user subscript throughout.

\runin{Raw sessions: the evidence that remains available}
A session $s$ is a timestamped sequence of user and agent turns, and
$\Hh=(s_1,\ldots,s_n)$ is the ordered session history. A closed session
is retained verbatim. This preserves details that no previous query
required and lets a later query revisit the original evidence. The
author of a turn is part of its provenance: an agent's unsupported
suggestion is not a fact about the user's activities or preferences.

\runin{First-order memory: reusable facts that queries reveal}
The fact memory $\Ff$ contains statements extracted from $\Hh$. A fact
has the minimal semantic form
\begin{equation}
\label{eq:new-fact}
f=(\mathrm{subject},\mathrm{predicate},\mathrm{object}).
\end{equation}
Each fact additionally has a stable identifier, source-span references,
and temporal qualifiers. Events receive their own identifiers, so
\texttt{(Alice, attended, hike-1)} and
\texttt{(hike-1, date, 1 March)} refer to the same event across sessions.
This representation supports properties and relations without
introducing a separate symbol family for each fact type.

Facts are acquired when a query needs them. An early query about an
event's organizer need not extract every event duration or every
preference mentioned in the history. A later query can reveal those
requirements and enrich $\Ff$. This use of the query to acquire durable
state is inspired by the query-driven evolution in EUR~\cite{eur}; here,
the evolving state consists of user facts and their use experiences.

\runin{Second-order memory: experience of using first-order memory}
The second-order memory $\Uu$ represents how first-order facts serve
this user's information needs. An experience distinguishes the facts
\emph{retrieved} from the facts \emph{used}, records the query context,
and associates execution outcomes and costs with that context. It can
therefore learn that the user's travel questions usually need event
dates and organizers, while preference questions use a different group
of facts. These experiences concern facts already in $\Ff$; they do not
add an unsupported belief about the user.

This distinction gives the architecture two forms of adaptation.
First-order memory improves \emph{what the agent knows about the user}.
Second-order memory improves \emph{what the memory system knows about
the user's memory needs}. A rarely retrieved fact can remain true, and
a frequently used fact can still require source verification. Usefulness
changes retrieval and computation priorities, while evidence determines
whether a statement can be returned as supported.

\begin{memoryexample}{Example 1: one history, three forms of state}
\textbf{Raw sessions $\Hh$.}
\par\texttt{S1, 2 March:} ``On 1 March I hiked outdoors on Pine Trail for four
hours. Ava organized the hike.''
\par\texttt{S2, 9 March:} ``On 8 March I took a two-hour indoor pottery
class at Clay Studio. Ben organized it.''
\par\texttt{S3, 16 March:} ``On 15 March I cycled outdoors around Lake
Park for three hours. Ava organized that ride too.''
\smallskip
\textbf{First query.} $q_1$: ``Who organized my Pine Trail hike?''
\par\textbf{New first-order facts.}
\par\texttt{f1: (Alice, attended, hike-1)}
\par\texttt{f2: (hike-1, location, Pine Trail)}
\par\texttt{f3: (hike-1, organizer, Ava)}
\par All three cite spans in \texttt{S1}. Durations are not needed yet.
\smallskip
\textbf{Retrieved memory.} $m_1$ returns \texttt{Ava}, supported by
\texttt{f1--f3}. The agent answers the user and reports using these facts.
\par\textbf{Second-order experience.} $\Uu$ remembers that this event
identity and organizer were used in an organizer query, together with
the observed verification outcome and execution cost.
\end{memoryexample}

\subsection{A Query Both Reads and Enriches Memory}
\label{subsec:new-evolution}

\runin{Demand, acquisition, and retention}
Let $\operatorname{Need}(q,\Ff)$ identify evidence obligations not yet
satisfied by the available facts. Obligations include missing properties,
missing relational links, unresolved versions, and unexamined parts of
a requested population. Query-driven acquisition is
\begin{align}
\label{eq:new-acquisition}
\Delta\Ff &= \operatorname{Extract}\bigl(\Hh,
                    \operatorname{Need}(q,\Ff)\bigr),\\
\Ff^{+} &= \operatorname{Merge}(\Ff,\Delta\Ff).
\label{eq:new-merge}
\end{align}
Extraction targets the unmet obligations and their source sessions.
Merge canonicalizes entity and event identities, verifies source spans,
normalizes dates and units, and retains reusable fact versions. Relative
dates are interpreted using the source session's timestamp. Provenance
and the source coverage of each extraction task are retained with the
facts, so an identical task over an unchanged source can be reused.

\runin{Versions preserve the user's history}
Two values with the same subject and predicate need not be inconsistent:
the user can visit multiple places or attend multiple events. Explicit
updates to a single-valued property, such as a current home address,
create temporally qualified versions. Historical queries select the
appropriate version. Ambiguous or conflicting statements remain
explicit until resolved; recency alone does not decide their truth.

\runin{The two updates}
For the $t$th memory query, acquisition and execution produce the
retrieved memory; the subsequent use report updates second-order memory:
\begin{align}
\label{eq:new-cycle}
(m_t,\Ff_t)&=\operatorname{Execute}
       (q_t,\Hh_t,\Ff_{t-1},\Uu_{t-1}),\\
\Uu_t&=\operatorname{Remember}
       (q_t,E_t^{r},E_t^{u},\ell_t,c_t;\Uu_{t-1}).
\label{eq:new-remember}
\end{align}
Here $E_t^{u}\subseteq E_t^{r}\subseteq\Ff_t$ are the facts used and
retrieved, respectively; $\ell_t$ is an observable evidence-risk
diagnostic and $c_t$ is measured memory-side execution cost. Used facts
include facts cited by the agent and facts contributing to a requested
computation. Until the agent supplies its use report, agent-level use
credit remains pending. A completed user--agent exchange is appended
to $\Hh$ when its session closes. A question can issue several memory
queries and therefore induce several updates before one answer.

\section{For Performance: Dependency-Aware Memory Querying}
\label{sec:new-performance}

Accurate memory access requires more than finding text similar to a
query. Some requests can be satisfied by selecting facts, some must
process a whole population, and others depend on the results of earlier
subrequests. We organize query execution by these \emph{evidence
dependencies}. Second-order experience guides which evidence to examine
first, while the dependency structure specifies what evidence must
ultimately be obtained and how it must be combined.

\subsection{Three Plans with a Defined Coverage Boundary}
\label{subsec:new-plans}

\runin{Retrieval, reduction, and composition}
The three plans differ in the computation required over evidence.
\emph{Evidence retrieval} (\ret) obtains facts satisfying a direct
information request. \emph{Evidence reduction} (\red) applies one
operation to a collection of evidence, such as counting, ranking,
comparing values, or summarizing events. \emph{Evidence composition}
(\comp) combines subrequests whose outputs have to be connected, such
as finding an event and then retrieving its organizer. Temporal and
preference queries use these same plans with different fact and time
constraints.

\begin{definition}[Memory query plans]
\label{def:new-plans}
A plan is a finite, acyclic expression generated by
\begin{equation}
\label{eq:new-plan-grammar}
\begin{split}
P\ ::={}&\ \ret(r)\\
       &\mid\ \red(g,P)\\
       &\mid\ \comp(g;P_1,\ldots,P_d),\quad d\geq2.
\end{split}
\end{equation}
Here $r$ is a scoped evidence request, and $g$ is an explicitly specified
reduction or composition operation. The output of a subplan can bind
an entity, event, or time in a later subplan. Each operation propagates
supporting facts and unresolved obligations with its result.
\end{definition}

Permitted reductions include cardinality, sums, extrema, ordering, and
source-backed summaries. Permitted compositions include identity-based
linking, temporal linking, and evidence-backed synthesis of multiple
subresults. A free-form generation prompt that hides an unspecified
multi-step computation does not qualify as a plan. In particular, the
plan must expose the facts and dependencies that its result needs.

\runin{Classification follows the outermost computation}
The planner parses $q$ into $P(q)$ and assigns its root:
\begin{equation}
\label{eq:new-route}
\rho(q)=
\begin{cases}
\ret, &P(q)=\ret(r),\\
\red, &P(q)=\red(g,P),\\
\comp,&P(q)=\comp(g;P_1,\ldots,P_d).
\end{cases}
\end{equation}
This makes the routing disjoint. A count over related records has a
\red{} root with a \comp{} child; a request for the organizer of the
longest event has a \comp{} root with a \red{} child. Nested plans are
executed recursively, rather than forced into one flat template.

\begin{proposition}[Structural coverage]
\label{prop:new-coverage}
Every plan in Definition~\ref{def:new-plans} has exactly one root class,
and can be evaluated using the three pipelines below, provided its
primitive requests and operations are supported.
\end{proposition}
\begin{proof}
The three constructors in Eq.~\eqref{eq:new-plan-grammar} are distinct,
so the outermost constructor determines one class. A retrieval is the
base case. A reduction evaluates its child and then its specified
operation. A composition evaluates its finite dependencies and then
combines their outputs. Structural induction on the finite plan gives
the result.
\end{proof}

The coverage claim concerns finite memory queries expressible through
these constructors. It does not assert that a parser understands every
natural-language utterance, that all evidence exists, or that an LLM
extracts it without error. Unsupported operations and unresolved
ambiguities yield a partial or unknown response. Requests for creative
content or external world knowledge remain the agent's responsibility.
The architecture thus has an explicit computational boundary instead
of relying on the apparent completeness of three keyword categories.

\begin{memoryexample}{Example 2: why the three plans cover different needs}
\textbf{Retrieve.} ``Who organized the Pine Trail hike?''
\par Select the organizer facts for the identified event.
\smallskip
\textbf{Reduce.} ``How many outdoor events did I attend in March?''
\par Retrieve the event population; evaluate the outdoor condition;
count distinct qualifying events.
\par ``Which lasted longer, the hike or the ride?'' uses the same plan
family with a maximum over two duration values.
\smallskip
\textbf{Compose.} ``Who organized the outdoor event that lasted longest?''
\par First reduce outdoor events by duration; then bind the winning
event in an organizer retrieval; compose the linked evidence.
\end{memoryexample}

\runin{A practical planner}
The agent supplies a structured request when possible. A deterministic
parser handles explicit attributes, scopes, and arithmetic operations;
ambiguous dependencies require a structured model-assisted parse.
Validation checks operation types, bindings, temporal scope, and
acyclicity before execution. A parse can be repaired within the query's
budget. Its LLM calls count as planning cost. Routing never uses a
benchmark category or a reference answer, and the revised method makes
no assumption that all planning is free.

\subsection{Second-Order Experience Guides Grounding}
\label{subsec:new-experience}

\runin{Remember exposure and use separately}
For a query context $z$, second-order memory maintains, for each fact,
an exposure count $n^r(f,z)$ and a credited-use count $n^u(f,z)$.
Contexts summarize the request's entities or topics, temporal intent,
and evidence dependencies. Related contexts share a documented
backoff rule; a new context begins with a neutral prior. The updates are
\begin{align}
\label{eq:new-use-update}
n^r(f,z)&\leftarrow n^r(f,z)+\mathbf{1}[f\in E^r],\\
n^u(f,z)&\leftarrow n^u(f,z)+w\,\mathbf{1}[f\in E^u],
\end{align}
where $0\leq w\leq1$ is the observed support credit for that use.
Credit requires a valid source link and a supported use; optional
explicit user feedback is stored separately. Samples agreeing on a
claim do not by themselves establish its truth. If $E^u$ is unavailable,
exposure can be recorded but use credit is not inferred from exposure.

\runin{Contextual usefulness}
With positive prior counts $a,b$, the estimated usefulness is
\begin{equation}
\label{eq:new-usefulness}
u(f\mid q,\Uu)=
\frac{a+n^u(f,z(q))}{a+b+n^r(f,z(q))}.
\end{equation}
This is a smoothed use ratio conditional on exposure, not a probability
that the fact is true. A fact repeatedly retrieved but never used loses
priority; a fact used in relevant contexts gains it. New facts retain
nonzero priority. Source validity and temporal eligibility are checked
before ranking.

\runin{Grounding combines present need and remembered use}
The grounding score is
\begin{equation}
\label{eq:new-ranking}
s(f\mid q)=\operatorname{rel}(f,q)\,
                 u(f\mid q,\Uu),
\end{equation}
where $\operatorname{rel}(f,q)\in[0,1]$ measures relevance to the current
request. The score orders eligible facts and the source sessions that
can extend them. A retained exploration portion of retrieval examines
new or weakly exposed facts: when available, each candidate batch
reserves one position for the least-exposed eligible fact with positive
current relevance, breaking ties by stable fact identifier. This
prevents earlier retrieval choices from permanently hiding alternatives.
Budget allocation can prioritize
promising candidates, but a population reduction cannot silently turn
the highest-ranked subset into the whole population.

The exposure/use distinction is essential: counting every keyword hit
as useful would reward the retrieval policy's own choices. Source
checks, use attribution, and exploration reduce this feedback bias;
they do not eliminate it. Evaluation must therefore include new-topic
queries and shuffled query orders, not only repeated queries.

\subsection{The Evidence Retrieval Pipeline}
\label{subsec:new-retrieve}

\runin{Scope, acquire, verify, return}
For $\ret(r)$, the system selects temporally eligible candidate facts
using Eq.~\eqref{eq:new-ranking}, determines unmet obligations, and
acquires missing facts using Eqs.~\eqref{eq:new-acquisition}--
\eqref{eq:new-merge}. It then verifies identity, source support, and
version consistency before returning
\begin{equation}
\label{eq:new-retrieve}
E=\operatorname{Ground}(r,\Hh,\Ff,\Uu),\qquad
v=\operatorname{Select}(r,E).
\end{equation}
Grounding can be repeated if the first source region is insufficient.
The response $m=(v,E,\sigma)$ includes the unresolved obligations and
keeps its source anchors. A supported value can be read directly from
a stored fact; no generation call is required merely to repeat it.
If candidate values conflict, the pipeline verifies the source or
returns the conflict for the agent to handle.

\begin{memoryexample}{Example 3: retrieval also acquires reusable knowledge}
\textbf{Request.} ``Who organized my Pine Trail hike?''
\par\textbf{Scope.} Identify the user's hike at Pine Trail.
\par\textbf{Acquire.} From \texttt{S1}, extract the attendance, location,
and organizer facts shown in Example 1.
\par\textbf{Verify.} The organizer span says ``Ava organized the hike'';
its event reference resolves to \texttt{hike-1}.
\par\textbf{Return.} $v=\texttt{Ava}$ with facts \texttt{f1--f3} and
\texttt{supported} status. The facts remain available for later queries.
\end{memoryexample}

\subsection{The Evidence Reduction Pipeline}
\label{subsec:new-reduce}

\runin{Enumerate, evaluate, reduce, verify coverage}
For $\red(g,P)$, the child plan first establishes the candidate evidence
collection. Records referring to the same event or entity are
canonicalized into a set of distinct candidates $D$. Each candidate
then receives the evidence and value needed by the reduction:
\begin{equation}
\label{eq:new-reduce}
x_e=\operatorname{Eval}(q,e,E),\qquad
v=g(\{(e,x_e):e\in D\}).
\end{equation}
Eval returns a source-supported value or an explicit unknown. For a
condition, $x_e\in\{1,0,\bot\}$ denotes supported yes, supported no, or
unresolved. For a comparison, $x_e$ is a normalized attribute value,
such as duration in hours. A semantic predicate can require an LLM;
sums, cardinalities, and extrema over verified values are deterministic.
The pipeline therefore exposes the uncertain judgments and performs
the final arithmetic explicitly.

\runin{A missing candidate is different from an unresolved condition}
An exact population result requires both candidate coverage and resolved
per-candidate evidence. Coverage is relative to the query's source
snapshot. A source-region inventory records which eligible sessions
have been examined for the relevant extraction obligations; new
sessions invalidate the corresponding coverage. In the absence of a
verified exclusion rule, all sessions in scope must be examined before
an exhaustive result is marked supported. Completing a top-$k$
retrieval is not a coverage certificate.

For a count over a certified candidate population $D$, the pipeline
can report an interval before every condition is resolved:
\begin{equation}
\label{eq:new-count-bounds}
\sum_{e\in D}\mathbf{1}[x_e=1]
\ \leq\ v\ \leq\
\sum_{e\in D}\mathbf{1}[x_e\neq0].
\end{equation}
If some source regions remain unexamined, even this upper bound is
not certified for the full population. The response reports a lower
bound and the open coverage obligation. A maximum similarly remains
partial when an unresolved candidate might exceed the current winner.
Source inspection and extraction are fallible in practice; coverage
metadata records the claimed scope and must be audited experimentally.

\begin{memoryexample}{Example 4: reduction discovers facts the first query missed}
\textbf{Request.} ``How many outdoor events did I attend in March?''
\par\textbf{Reuse.} The attendance and identity of \texttt{hike-1} are
already in $\Ff$. Its duration is unnecessary for this query.
\par\textbf{Extend.} Examine \texttt{S1--S3}; acquire dates, activity
types, and the attendance facts for \texttt{class-1} and \texttt{ride-1}.
\par\textbf{Evaluate each distinct event.}
\par\texttt{hike-1: hiking at Pine Trail; outdoor = yes.}
\par\texttt{class-1: indoor pottery; outdoor = no.}
\par\texttt{ride-1: cycling outdoors; outdoor = yes.}
\par\textbf{Reduce.} $v=1+0+1=2$, with an explicit three-event witness.
\par\textbf{Remember use.} The reduction used both positive evidence
and the pottery exclusion evidence. All contribute use experience;
the answer text alone would hide the role of the excluded candidate.
\end{memoryexample}

\subsection{The Evidence Composition Pipeline}
\label{subsec:new-compose}

\runin{Decompose, bind, execute, connect}
A composition exposes a finite dependency graph of subplans. In
topological order, earlier outputs bind identifiers or constraints in
later requests. Its result is
\begin{align}
\label{eq:new-compose}
(m_j,\Ff)&=\operatorname{Execute}
       (P_j[\mathrm{bindings}],\Hh,\Ff,\Uu),\\
v&=g(v_1,\ldots,v_d),\qquad E=\bigcup_{j=1}^{d}E_j.
\end{align}
The union is a candidate support set. The pipeline checks each link
and retains the facts actually supporting the composition; shared
facts and shared subplans are deduplicated. One subrequest can retrieve
events, another can compare their durations, and a later one can
retrieve a property of the selected event. Evidence from different
sessions is connected through stable identities and temporal
qualifiers, rather than by concatenating passages and assuming the
model aligns them.

A composition is supported only when all required links and required
child outputs are supported. If the longest event is unresolved, its
organizer cannot be returned as the organizer of the longest event,
even when a retrieved organizer fact itself is valid. Optional evidence
can remain partial without blocking a result whose specified operation
does not depend on it. This distinction is part of the plan.

\begin{memoryexample}{Example 5: the result of one request becomes the next scope}
\textbf{Request.} ``Who organized the outdoor event that lasted longest?''
\par\textbf{Subplan 1.} Reuse the outdoor-event population from Example 4.
Extract the missing durations from \texttt{S1} and \texttt{S3}:
\par\texttt{(hike-1, duration-hours, 4)}
\par\texttt{(ride-1, duration-hours, 3)}
\par Reduce by maximum; the winner is \texttt{hike-1}.
\par\textbf{Subplan 2.} Bind \texttt{event = hike-1} in the organizer
request. Reuse \texttt{f3: (hike-1, organizer, Ava)}.
\par\textbf{Connect.} The winner identifier and organizer subject match.
Return \texttt{Ava}, supported by the population, the durations, and
the organizer link. The reply needs no new organizer extraction.
\end{memoryexample}

\runin{What experience changes, and what the plan preserves}
Second-order memory affects the order and allocation of evidence work
in all three pipelines. It cannot remove a required population member,
change a reduction predicate, or replace a dependency link. Under
complete evidence and exact primitive operations, changing retrieval
order leaves the result unchanged. Under a finite budget, the system
may return partial evidence sooner, and its stated completeness changes
accordingly. Neither result invariance under truncation nor automatic
accuracy improvement is assumed.

\section{For Efficiency: Experience-Guided Cost-Aware Execution}
\label{sec:new-efficiency}

The plan specifies what computation a memory query requires; it does
not determine how many LLM executions that particular query needs.
A reduction over cached numerical facts can be cheaper than a retrieval
with ambiguous identity. The executor therefore budgets the query
using its current evidence and second-order experience, and revises
the budget after each execution.

\subsection{Budget the Query, Not the Plan Label}
\label{subsec:new-budget}

\runin{Cost includes acquiring the missing memory}
For $k$ stochastic execution rounds, memory-side cost is
\begin{equation}
\label{eq:new-cost}
C(q,k)=C_{\mathrm{plan}}(q)+C_{\mathrm{acquire}}(q)
              +\sum_{j=1}^{k}C_{\mathrm{exec}}(q,j).
\end{equation}
Planning, extraction, verification, and model-assisted predicate
evaluation are counted wherever they occur. Deterministic selection
and arithmetic do not add an LLM call. Cached facts and unchanged
coverage are reused across rounds. If new evidence is acquired during
execution, its cost is added to $C_{\mathrm{acquire}}$, and the affected
execution state is revised. Cost can be measured in LLM calls or
priced tokens, but one chosen unit must be used consistently in the
budget objective. Latency is measured separately. The agent's final
answer generation is reported as an additional end-to-end cost.

\runin{Query features are observed before allocating executions}
The context $z(q)$ includes candidate-set size, dependency depth,
unresolved evidence, conflicting versions, and the fraction of usable
facts already cached. It also includes the query topic used for
second-order retrieval. These features distinguish queries within
the same plan family. The executor reads experience for matching
contexts to estimate how risk and cost change with additional work.

\runin{Risk is observable evidence risk}
Let $\ell(q,m)\in[0,1]$ measure an execution's unresolved evidence risk.
The diagnostic combines unfulfilled requirements, unsupported result
claims, unresolved dependencies, and disagreement on required semantic
judgments. Specifically,
\begin{equation}
\label{eq:new-risk}
\ell(q,m)=\max\{\operatorname{miss}(q,E),\,
  \operatorname{unsup}(v,E),\,\operatorname{dis}(q)\}.
\end{equation}
Here $\operatorname{miss}$ is the fraction of mandatory evidence
obligations still open, including required population coverage and
dependency links; $\operatorname{unsup}$ is the fraction of returned
atomic claims without checked support; and $\operatorname{dis}$ is the
fraction of required semantic decisions with conflicting observed
judgments. A component with no applicable obligations or claims is
zero. An unexecuted mandatory decision remains a missing obligation,
even though it has not yet produced disagreement. The checks are fixed
before evaluation and their components are reported separately.
An exact deterministic reduction with verified
inputs can have zero diagnostic risk without repeated generation.
Two ungrounded samples giving the same response still fail the source
support checks. This diagnostic is an optimization surrogate; it is
not an oracle for answer correctness.

\runin{An experience-conditioned execution budget}
Given predicted risk $\widehat L(q,k\mid\Uu)$ and cost
$\widehat C(q,k\mid\Uu)$, the initial budget is
\begin{equation}
\label{eq:new-adaptive-budget}
K(q)=\argmin_{0\leq k\leq K_{\max}(q)}
\left\{\widehat L(q,k\mid\Uu)
        +\lambda\widehat C(q,k\mid\Uu)\right\}.
\end{equation}
The nonnegative multiplier $\lambda$ controls the tradeoff between
evidence risk and measured cost. $K_{\max}(q)$ is a resource cap, not a
fixed budget assigned to a plan class. $k=0$ is feasible when cached
evidence and deterministic operations can already produce a supported
response. If a query needs extraction, that acquisition cost remains
part of the objective even when no subsequent sampling is needed.
In a discrete tie, the smaller budget is chosen.

Risk and cost estimates are computed from observed outcomes of
similar contexts. Sparse contexts back off to a declared prior, and
estimates retain uncertainty intervals. Trials not executed have no
observed outcomes; the executor does not manufacture counterfactual
risk reductions. The budget is consequently a decision under estimated
uncertainty, not a claim to know the optimal number of calls in advance.

\subsection{Revise the Budget by the Value of More Computation}
\label{subsec:new-value}

\runin{Spend on the unresolved part of the pipeline}
After a round, the executor checks evidence obligations and disputed
judgments. It can acquire an unexamined source region or rerun an
uncertain semantic block. Verified fact lookups and deterministic
reductions are reused. If acquired evidence changes an input, dependent
blocks are invalidated and recomputed. This avoids spending another
full pass merely because one event's condition is ambiguous.

\runin{A sequential decision}
Let $a$ be an admissible next action, and let
$\widehat{\Delta L}_k(a)$ be its predicted reduction in the current
evidence-risk diagnostic, conditional on all work so far. The decision is
\begin{equation}
\label{eq:new-value}
\begin{split}
a_k^{*}&=\argmax_a
     \left\{\widehat{\Delta L}_k(a)-\lambda\widehat c_k(a)\right\},\\
\mathrm{continue}&\quad\text{iff}\quad
\max_a\left\{\widehat{\Delta L}_k(a)
                       -\lambda\widehat c_k(a)\right\}>0.
\end{split}
\end{equation}
An action is admissible only if it respects the remaining resource cap,
the query's source scope, and the plan's dependencies. On stopping,
the current response is returned with its actual support status. An
unmet obligation never becomes satisfied solely because the resource
cap was reached.

Equation~\eqref{eq:new-value} implements a receding-horizon,
one-step value-of-computation policy. The initial budget in
Eq.~\eqref{eq:new-adaptive-budget} is its lookahead envelope and is
re-estimated after observations. When a sequence of actions is needed
before any improvement is visible, that sequence is evaluated as a
bounded admissible bundle; otherwise a one-step rule could stop before
completing a useful dependency. Global optimality is not claimed for
the practical estimator.

\runin{Agreement is specific to a supported computation}
For repeated semantic judgments, outputs are normalized to the plan's
required result: event labels for a predicate, normalized values for
a comparison, or canonical facts for a synthesis. The executor measures
disagreement on those values, while checking source support and
coverage independently. It never compares entire explanatory paragraphs
as a proxy for factual equivalence. Agreement can lower the predicted
value of another execution, but its stopping effect depends on the
query's unresolved obligations and observed cost.

\begin{proposition}[The local stopping condition]
\label{prop:new-stopping}
If predicted risk changes and costs equal their conditional expectations,
and an admissible action is followed by stopping, taking the action
decreases the expected risk-plus-cost objective exactly when
$\widehat{\Delta L}_k(a)>\lambda\widehat c_k(a)$.
\end{proposition}
\begin{proof}
Let the current diagnostic be $L$. Stopping incurs $L$. Taking $a$
and then stopping incurs
$L-\widehat{\Delta L}_k(a)+\lambda\widehat c_k(a)$ in expectation.
Their difference gives the stated condition. Past cost is already
incurred and cancels from this comparison.
\end{proof}
The proposition justifies the decision rule for its specified local
objective. It does not establish that a diagnostic predicts all
answer errors, or that estimated action values are exact.

\subsection{Execution Experience Becomes Second-Order Memory}
\label{subsec:new-cost-experience}

\runin{Remember what an execution bought}
In addition to fact-use counts, $\Uu$ stores execution summaries for
context and action pairs: number of observations, mean observed risk
change, mean measured cost, and their dispersion. For any recorded
outcome $x$, its mean is updated incrementally:
\begin{equation}
\label{eq:new-mean-update}
n\leftarrow n+1,\qquad
\overline{x}\leftarrow\overline{x}
                    +\frac{x-\overline{x}}{n}.
\end{equation}
The same update applies to cost and observed risk changes. A prefix
outcome is associated only with the work actually performed. The
observed action's source obligations and cache state are part of its
context, so costs from a cold query do not masquerade as costs of a
warm query. Fact-use summaries guide grounding; execution summaries
guide the next budget. Both are memories of serving the same user.

\runin{Shared evidence does not create independent votes}
Repeated evaluations can share a frozen fact snapshot while drawing
new semantic judgments. Those judgments provide conditional
disagreement evidence, not statistically independent verification of
the underlying extraction. If a source interpretation is in doubt,
the executor must revisit that source instead of repeatedly sampling
from an erroneous cached fact. Successful acquisition enriches $\Ff$;
its verified downstream use and measured cost enrich $\Uu$.

\begin{memoryexample}{Example 6: two reductions receive different budgets}
\textbf{Warm comparison.} The durations \texttt{4 hours} and
\texttt{3 hours} from Example 5 are cached and source-verified.
``Which event lasted longer?'' needs a deterministic maximum. Its
additional stochastic execution budget can be $K(q)=0$.
\smallskip
\textbf{Uncertain count.} A later session describes another activity
whose outdoor status is unclear. ``How many outdoor events?'' now has
a coverage update and an unresolved candidate condition. The executor
first pays for the missing evidence. If one semantic evaluation leaves
the condition disputed and second-order experience predicts useful
risk reduction, it allocates another execution to that condition.
If the source cannot resolve it, the count remains partial.
\smallskip
Both requests have a \red{} root. Their allocation follows evidence,
uncertainty, reuse, and predicted cost rather than a class-specific
constant. These budgets illustrate the decision rule; they are not
reported experimental measurements.
\end{memoryexample}

\runin{End-to-end behavior}
A query identifies evidence needs, acquires reusable facts, and executes
a dependency-specific pipeline. The executor allocates further work
where its predicted benefit warrants its cost. The agent receives a
supported or explicitly incomplete memory response and generates the
user-facing answer. The fact-use and execution experiences are then
remembered, and the exchange joins the raw history. Later queries can
reuse both the acquired knowledge and the experience of acquiring and
using it.

\section{Experiments and Evaluation}
\label{sec:new-experiments}

Our evaluation separates the available empirical evidence from the
experiments needed for the revised architecture. The released artifact
contains an earlier prototype and four sets of per-question evaluation
records. We recompute their aggregates and inspect their category
distributions. This audit is evidence about that prototype, while the
fact-level use experience, recursive composition, and dynamic budget
specified above require a matched implementation and evaluation.

\subsection{Artifact, Metrics, and Reproducibility}
\label{subsec:new-artifact}

\runin{Available records}
The archived runs contain 1,540 LoCoMo records, 451 records under the
artifact's LongMemEval-v label, 30 EvoMem records, and 240 DynamicMem
records. Each has a question identifier, a category, a generated answer,
and a score in $[0,1]$. LoCoMo, EvoMem, and DynamicMem also have a
paired Full-Context score. LongMemEval-v does not. We use the artifact's
variant name and category identifiers without assuming that its
451-record subset is the standard LongMemEval benchmark configuration.

\runin{Aggregation}
For a set of $N$ records with scores $s_i$, the reported score is
\begin{equation}
\label{eq:new-score}
\operatorname{Score}=\frac{100}{N}\sum_{i=1}^{N}s_i.
\end{equation}
The archive describes token F1 for LoCoMo, phrase-set F1 for
LongMemEval-v, and loose F1 for EvoMem and DynamicMem. These scales are
not averaged across benchmarks. We verify record counts, unique
identifiers, score ranges, and aggregate calculations. The archive does
not include the source questions, reference answers, complete scoring
scripts, or run configuration; consequently, this audit does not
independently re-score the generated answers or establish the exact
backbone and prompt conditions of every run.

\runin{The released implementation}
The source archive provides cached fact-triple extraction, per-session
entity matching, query-category-dependent rendering switches, a
session-use counter, and lexical query routing. Its demonstration runs
over eight synthetic sessions. The demonstration without a configured
LLM exercises routing, prompt construction, and experience persistence,
but produces no extracted triples or answered-query accuracy. The
archive does not contain the revised fact-use reporting interface,
recursive dependency executor, or query-specific risk--cost allocator.

\begin{table*}[t]
\centering\small
\caption{Recomputed scores of the archived prototype. Full-Context is
the paired baseline present in the records. Differences are percentage
points on each benchmark's own metric. These runs predate the revised
mechanisms; no dynamic-budget gain is inferred from them.}
\label{tab:new-archived-results}
\begin{tabular}{@{}lrlrrr@{}}
\toprule
Archived benchmark & $N$ & Metric & Full-Context & Prototype & Difference\\
\midrule
LoCoMo & 1,540 & Token F1 & 52.24 & 58.76 & $+6.53$\\
LongMemEval-v & 451 & Phrase-set F1 & -- & 10.31 & --\\
EvoMem & 30 & Loose F1 & 16.68 & 24.60 & $+7.92$\\
DynamicMem & 240 & Loose F1 & 33.53 & 34.93 & $+1.40$\\
\bottomrule
\end{tabular}
\end{table*}

\subsection{What the Archived Runs Establish}
\label{subsec:new-archive-results}

Table~\ref{tab:new-archived-results} reproduces the archived means to
two decimal places. On the three runs with paired scores, the prototype
exceeds Full-Context by 6.53, 7.92, and 1.40 points, respectively.
The paired record-level win/tie/loss counts are 630/510/400 for LoCoMo,
28/0/2 for EvoMem, and 108/98/34 for DynamicMem. These counts describe
the archived answer scores; repeated questions from one user are not
treated as independent evidence for a statistical significance claim.

\runin{The aggregate hides different behaviors}
Table~\ref{tab:new-category-results} retains the archive's category
identifiers. In LoCoMo, the prototype improves on \texttt{cat=1} and
\texttt{cat=4}, but falls below Full-Context on \texttt{cat=2} and
\texttt{cat=3}. All three DynamicMem category means improve, with
substantially different absolute score levels. The original category
labels do not establish the new plan root: a new dependency annotation
is needed before any Retrieve/Reduce/Compose decomposition is reported.

\begin{table}[t]
\centering\small
\caption{Archived category scores, recomputed without remapping them to
the revised plan classes. FC denotes Full-Context.}
\label{tab:new-category-results}
\begin{tabular}{@{}lrrr@{}}
\toprule
Category & $N$ & FC & Prototype\\
\midrule
\multicolumn{4}{@{}l}{\emph{LoCoMo}}\\
\texttt{cat=1} & 282 & 40.43 & 50.84\\
\texttt{cat=2} & 321 & 52.10 & 50.52\\
\texttt{cat=3} & 96 & 35.93 & 32.60\\
\texttt{cat=4} & 841 & 58.11 & 67.55\\
\midrule
\multicolumn{4}{@{}l}{\emph{DynamicMem}}\\
Habits & 79 & 65.78 & 66.96\\
Preferences & 57 & 20.19 & 21.97\\
User attributes & 104 & 16.34 & 17.70\\
\bottomrule
\end{tabular}
\end{table}

\runin{Observed failures need their own analysis}
Fifty-nine LoCoMo records have empty generated answers. In the
LongMemEval-v archive, all 157 records in the three
\texttt{*-abs} categories and \texttt{errors-gotchas} have zero scores.
Without gold answers and execution traces, these observations cannot
be attributed specifically to retrieval, abstention, extraction, or
answer formatting. The revised response interface makes these failure
types explicit and motivates preserving source and execution traces
alongside final answer scores.

\subsection{Matched Evaluation of the Revised System}
\label{subsec:new-evaluation-protocol}

The revised system requires three groups of experiments, corresponding
to architecture, performance, and efficiency. These are the protocol
for the next implementation rather than completed results.

\runin{Architecture: does query-driven memory acquire reusable knowledge?}
Compare raw-session retrieval, upfront fact extraction, and query-driven
fact acquisition under the same agent and evidence scope. Report answer
quality, fact-source precision, recall of query-required facts, and
the rate at which later queries reuse earlier fact acquisitions.
Separate cold queries from warm queries. Preserve historical versions
and test both current and historical requests. A fixed summary should
not receive a larger inference budget merely to compensate for omitted
facts without recording that additional cost.

\runin{Performance: do dependencies and use experience improve evidence?}
Compare a generic retrieve-then-read pipeline with the three recursive
pipelines, keeping the fact state, agent, and execution cap matched.
Annotate a held-out sample with scopes, root classes, candidate sets,
and dependency links. Report root classification, missing dependencies,
source coverage, and final answer quality. Ablate fact-use experience,
exposure/use separation, and exploration. Evaluate real query order,
shuffled order, and new-topic queries with an empty or warm $\Uu$.
Hold the first-order cache fixed in one comparison so that improved
fact coverage is not mistaken for improved second-order ranking.

\runin{Efficiency: does the individual-query allocator spend better?}
Compare uniform execution counts, fixed class budgets, and
Eq.~\eqref{eq:new-adaptive-budget}. Use identical evidence acquisition
and answer generation. Plot quality against measured cost across
resource caps and $\lambda$ values, including acquisition, planning,
verification, and the final agent call. Evaluate full reruns against
selective block re-execution. Report the calibration of predicted risk
changes against their observed values and the relationship between the
diagnostic and actual correctness. In an offline audit subset, execute
multiple budget prefixes to obtain genuine additional-work outcomes;
only the prefixes actually observed can train or assess the estimator.

\runin{Statistical and resource controls}
Freeze benchmark splits, model versions, prompts, category annotation,
and cost units before comparison. Keep user partitions isolated.
Report uncertainty by resampling users or conversation groups when
those identifiers are available. Model-call counts, input/output
tokens, wall-clock latency, and persistent state size are distinct
measurements. Session-level extraction cost must be allocated to the
workload or reported separately, and never omitted only for one system.
Memory capacity should be varied independently from the number of
queries, because first-order fact growth and second-order experience
growth need not have the same scaling behavior.

\subsection{Limits and Error Analysis}
\label{subsec:new-limitations}

The model exposes several limits that evaluation must measure.
Extraction can miss or misinterpret facts; stable identity resolution
can incorrectly merge events; and a semantic predicate can remain
ambiguous even after repeated execution. Second-order ranking can
reinforce an early exposure bias, especially under a small retrieval
budget. Dynamic allocation can stop too early when its risk estimator
is poorly calibrated or its action horizon is too short. A complete
population request can require inspection proportional to the in-scope
history: call caps do not make that source work constant.

These limits distinguish structural claims from empirical ones.
Proposition~\ref{prop:new-coverage} establishes coverage of the declared
plan language, not natural-language parsing accuracy.
Proposition~\ref{prop:new-stopping} establishes a conditional local
decision rule, not a bound on answer error. The available archive
cannot validate fact-use attribution, query-position adaptation,
dynamic-budget savings, or revised scalability. Those claims require
the traces and matched experiments described above.

\section{Related Work}
\label{sec:new-related}

\runin{Conversation and fact memory}
MemGPT manages information across context and external memory
tiers~\cite{memgpt}. Mem0 builds persistent memory by extracting and
updating salient conversational facts~\cite{mem0}. Zep represents
time-dependent relationships through a temporal knowledge
graph~\cite{zep}, while A-MEM builds linked notes that evolve with
new experiences~\cite{amem}. These systems provide important fact
and context representations. \sys{} focuses on the relationship
between query-induced first-order facts and second-order experience
of using those facts. Its second-order state refers to evidence
exposure, downstream use, and execution cost, rather than simply
adding another semantic description of the conversation.

\runin{Retrieval control and memory queries}
LongMemEval separates memory design into indexing, retrieval, and
reading, and evaluates information extraction, multi-session reasoning,
temporal reasoning, knowledge updates, and abstention~\cite{longmemeval}.
MemR$^3$ uses an evidence-gap tracker and a controller that selects
retrieval, reflection, or answering~\cite{memr3}. These are relevant
precedents for explicit memory access and iterative evidence gathering.
The three classes in this paper instead describe the structure of the
requested evidence computation: selecting facts, reducing a collection,
or composing dependencies. They are plan constructors, not the
controller's possible next actions. The cost controller acts on these
plans and retains user-specific experience of its actions.

\runin{Query-driven acquisition}
EUR~\cite{eur} argues that incoming queries can reveal missing concepts
and facts and leave reusable state for future workloads. We borrow
the query-driven acquisition principle and the organization of a
running example through successive state changes. Our interface serves
an agent's information needs over conversational sessions: first-order
state contains user facts, and second-order state remembers their use.
The contribution under study is this coupled memory behavior, rather
than treating conversation history as a schema-first store.

\runin{Repeated reasoning and computational budgets}
Self-consistency samples multiple reasoning paths and aggregates their
answers~\cite{selfconsistency}. Agreement is useful evidence about
variation among outputs, but cannot verify a shared false premise.
Our execution objective conditions additional work on source support,
evidence coverage, unresolved dependencies, and cost experience.
A query can reuse verified deterministic results without resampling,
or spend on one unresolved predicate or source region. Empirical
comparison with fixed and adaptive sampling policies is necessary to
establish whether the additional state and decisions are worthwhile.

\runin{Evaluation of long-term agent memory}
LoCoMo studies long-term conversational recall and related tasks across
many sessions~\cite{locomo}. LongMemEval emphasizes multiple forms of
memory reasoning~\cite{longmemeval}, and MemBench broadens memory
evaluation across scenarios and question types~\cite{membench}.
Their category labels are not interchangeable with the plan roots
defined here. A query-level dependency annotation is needed to
evaluate structural routing, while chronological and shuffled
workloads are needed to evaluate second-order adaptation.

\section{Conclusion}
\label{sec:new-conclusion}

\sys{} organizes agent memory through raw conversations, query-driven
facts, and experience of using those facts. This separation preserves
the evidence for future needs, lets first-order knowledge grow when a
query requires it, and lets second-order knowledge guide later memory
access. Three recursive evidence pipelines expose direct requests,
population computations, and linked dependencies. A cost-aware executor
allocates work to the individual query and updates its allocation from
observed evidence and execution experience.

The resulting model has explicit interfaces, state transitions,
support obligations, and a computational coverage boundary. The
archived prototype records establish a reproducible empirical starting
point, including category-level failures. Implementing the revised
fact-use interface and dynamic executor, then evaluating them under
matched source, model, and cost conditions, is the remaining step for
validating the proposed performance and efficiency benefits.

\begin{thebibliography}{99}
\bibitem{mem0}
Prateek Chhikara, Dev Khant, Saket Aryan, Taranjeet Singh, and Deshraj
Yadav. 2025. Mem0: Building Production-Ready AI Agents with Scalable
Long-Term Memory. \emph{arXiv:2504.19413}.
\url{https://arxiv.org/abs/2504.19413}.

\bibitem{memr3}
Xingbo Du, Loka Li, Duzhen Zhang, and Le Song. 2025.
MemR$^3$: Memory Retrieval via Reflective Reasoning for LLM Agents.
\emph{arXiv:2512.20237}.
\url{https://arxiv.org/abs/2512.20237}.

\bibitem{eur}
Nan Tang and Haixun Wang.
Back to the Future: The Evolving Universal Relation Model.
Undated manuscript supplied as a reference.

\bibitem{memgpt}
Charles Packer, Sarah Wooders, Kevin Lin, Vivian Fang, Shishir G. Patil,
Ion Stoica, and Joseph E. Gonzalez. 2023.
MemGPT: Towards LLMs as Operating Systems. \emph{arXiv:2310.08560}.
\url{https://arxiv.org/abs/2310.08560}.

\bibitem{zep}
Preston Rasmussen, Pavlo Paliychuk, Travis Beauvais, Jack Ryan, and
Daniel Chalef. 2025.
Zep: A Temporal Knowledge Graph Architecture for Agent Memory.
\emph{arXiv:2501.13956}. \url{https://arxiv.org/abs/2501.13956}.

\bibitem{amem}
Wujiang Xu, Zujie Liang, Kai Mei, Hang Gao, Juntao Tan, and Yongfeng
Zhang. 2025. A-MEM: Agentic Memory for LLM Agents.
\emph{arXiv:2502.12110}. \url{https://arxiv.org/abs/2502.12110}.

\bibitem{selfconsistency}
Xuezhi Wang, Jason Wei, Dale Schuurmans, Quoc Le, Ed Chi, Sharan
Narang, Aakanksha Chowdhery, and Denny Zhou. 2022.
Self-Consistency Improves Chain of Thought Reasoning in Language Models.
\emph{arXiv:2203.11171}. \url{https://arxiv.org/abs/2203.11171}.

\bibitem{locomo}
Adyasha Maharana, Dong-Ho Lee, Sergey Tulyakov, Mohit Bansal, Francesco
Barbieri, and Yuwei Fang. 2024. Evaluating Very Long-Term Conversational
Memory of LLM Agents. \emph{ACL}.
\url{https://arxiv.org/abs/2402.17753}.

\bibitem{longmemeval}
Di Wu, Hongwei Wang, Wenhao Yu, Yuwei Zhang, Kai-Wei Chang, and Dong Yu.
2025. LongMemEval: Benchmarking Chat Assistants on Long-Term Interactive
Memory. \emph{ICLR}. \url{https://arxiv.org/abs/2410.10813}.

\bibitem{membench}
Haoran Tan, Zeyu Zhang, Chen Ma, Xu Chen, Quanyu Dai, and Zhenhua Dong.
2025. MemBench: Towards More Comprehensive Evaluation on the Memory of
LLM-based Agents. \emph{Findings of ACL}.
\url{https://arxiv.org/abs/2506.21605}.
\end{thebibliography}

\end{document}
